\documentclass[tikz,border=8pt]{standalone}
\usepackage{tikz}
\usepackage{amsmath}
\usepackage{amssymb}
\usepackage{pgfplots}
\pgfplotsset{compat=1.18}
\usepackage{ctex}
\usetikzlibrary{calc,positioning,arrows.meta,matrix,trees}

% LC 01-Matrix / BFS 最短路径
% 4x4 网格，起点 (0,0) 绿色，终点 (3,3) 橙色
% BFS 层级用浅蓝渐变：层1=blue!15，层2=blue!30，层3=blue!45，层4=blue!60

\begin{document}
\begin{tikzpicture}[
  every node/.style={font=\small},
  gcell/.style={draw=gray!60, rectangle, minimum size=0.9cm, thick, fill=white},
  startcell/.style={draw=green!60!black, rectangle, minimum size=0.9cm, very thick, fill=green!20},
  endcell/.style={draw=orange!70!black, rectangle, minimum size=0.9cm, very thick, fill=orange!20},
  L1/.style={draw=blue!60, rectangle, minimum size=0.9cm, thick, fill=blue!15},
  L2/.style={draw=blue!70, rectangle, minimum size=0.9cm, thick, fill=blue!25},
  L3/.style={draw=blue!80, rectangle, minimum size=0.9cm, thick, fill=blue!35},
  L4/.style={draw=blue!90, rectangle, minimum size=0.9cm, thick, fill=blue!50},
  wallarr/.style={draw=gray!40, rectangle, minimum size=0.9cm, thick, fill=gray!20},
  ptr/.style={->,>=Stealth,thick},
]

%% 标题
\node[font=\large\bfseries] at (4.0, 1.5)
  {BFS 网格最短路径（01-Matrix / Word Ladder）};
\node[font=\small,gray] at (4.0, 0.9)
  {BFS 逐层扩展，数字=到起点最短距离；绿=起点，橙=终点};

%% 4x4 网格（行0在顶，col0在左）
%% 障碍格用灰色，其余按 BFS 层染色
%% 层0: (0,0); 层1: (0,1),(1,0); 层2: (0,2),(1,1),(2,0);
%% 层3: (0,3),(1,2),(2,1),(3,0); 层4: (1,3),(2,2),(3,1); 层5: (2,3),(3,2); 层6: (3,3)

\def\cs{1.0} % cell spacing

% 定义每格样式和标注值
% (row, col) -> style, label
% Row 0
\node[startcell]    at (0*\cs, 0)   {\textbf{0}};
\node[L1]           at (1*\cs, 0)   {1};
\node[L2]           at (2*\cs, 0)   {2};
\node[L3]           at (3*\cs, 0)   {3};
% Row 1
\node[L1]           at (0*\cs, -1*\cs)  {1};
\node[L2]           at (1*\cs, -1*\cs)  {2};
\node[L3]           at (2*\cs, -1*\cs)  {3};
\node[L4]           at (3*\cs, -1*\cs)  {4};
% Row 2
\node[L2]           at (0*\cs, -2*\cs)  {2};
\node[L3]           at (1*\cs, -2*\cs)  {3};
\node[L4]           at (2*\cs, -2*\cs)  {4};
\node[L4,draw=blue!90,fill=blue!45] at (3*\cs, -2*\cs)  {5};
% Row 3
\node[L3]           at (0*\cs, -3*\cs)  {3};
\node[L4]           at (1*\cs, -3*\cs)  {4};
\node[L4,draw=blue!90,fill=blue!45] at (2*\cs, -3*\cs)  {5};
\node[endcell]      at (3*\cs, -3*\cs)  {\textbf{6}};

%% 行列标注
\foreach \c/\lb in {0/0, 1/1, 2/2, 3/3} {
  \node[font=\tiny,gray] at (\c*\cs, 0.62) {col\c};
}
\foreach \r/\lb in {0/0, 1/1, 2/2, 3/3} {
  \node[font=\tiny,gray,anchor=east] at (-0.55, -\r*\cs) {row\r};
}

%% BFS 路径箭头（从 (0,0) 到 (3,3) 的一条最短路径）
\def\px{0.45}  % half cell offset for arrow anchors
\draw[ptr,red!80!black,very thick]
  (0*\cs+\px, 0) -- (1*\cs-\px, 0);
\draw[ptr,red!80!black,very thick]
  (1*\cs, -\px) -- (1*\cs, -1*\cs+\px);
\draw[ptr,red!80!black,very thick]
  (1*\cs+\px, -1*\cs) -- (2*\cs-\px, -1*\cs);
\draw[ptr,red!80!black,very thick]
  (2*\cs, -1*\cs-\px) -- (2*\cs, -2*\cs+\px);
\draw[ptr,red!80!black,very thick]
  (2*\cs+\px, -2*\cs) -- (3*\cs-\px, -2*\cs);
\draw[ptr,red!80!black,very thick]
  (3*\cs, -2*\cs-\px) -- (3*\cs, -3*\cs+\px);

%% 图例（右侧）
\node[font=\small\bfseries] at (7.0, 0) {图例};
\node[startcell,minimum size=0.6cm,font=\scriptsize] at (6.3, -0.7) {};
\node[font=\scriptsize,anchor=west] at (6.8, -0.7) {起点 (0,0)};
\node[endcell,minimum size=0.6cm,font=\scriptsize]   at (6.3, -1.3) {};
\node[font=\scriptsize,anchor=west] at (6.8, -1.3) {终点 (3,3)};
\node[L1,minimum size=0.6cm]  at (6.3, -1.9) {};
\node[font=\scriptsize,anchor=west] at (6.8, -1.9) {BFS 层 1};
\node[L2,minimum size=0.6cm]  at (6.3, -2.5) {};
\node[font=\scriptsize,anchor=west] at (6.8, -2.5) {BFS 层 2};
\node[L4,minimum size=0.6cm]  at (6.3, -3.1) {};
\node[font=\scriptsize,anchor=west] at (6.8, -3.1) {BFS 层 3+};

%% BFS 队列演化（下方文字描述）
\node[font=\small\bfseries] at (4.0, -4.5) {BFS 队列演化};

\node[draw=gray!50,fill=white,rounded corners=3pt,font=\scriptsize,
      inner sep=6pt,align=left]
  at (4.0, -5.9)
  {\textbf{层0} 初始队列：\{(0,0)\}，距离=0\\[3pt]
   \textbf{层1} 出队(0,0)，入队相邻：\{(0,1),(1,0)\}\\[3pt]
   \textbf{层2} 出队(0,1),(1,0)，入队未访问相邻：\{(0,2),(1,1),(2,0)\}\\[3pt]
   \textbf{层3} 继续扩展…每个格记录到起点距离\\[3pt]
   \textbf{终止} 当终点(3,3)出队，最短路径长度=6};

%% 算法说明
\node[draw=gray!50,fill=white,rounded corners=3pt,font=\small,
      inner sep=8pt,align=left]
  at (4.0, -8.0)
  {\textbf{BFS 最短路径 O(m×n) 时间空间：}\\[4pt]
   将起点加入队列，distance 设为0；\\[2pt]
   每次出队一格，将未访问的四邻格入队，距离+1；\\[2pt]
   BFS 保证首次到达某格时路径最短（无权图）};

\end{tikzpicture}
\end{document}
